% SHARD-ID: CA-54-QUBIT-HAMILTONIAN-FAMILY-CONVENTIONS
% SHARD-TITLE: Pauli Coefficients for Scan Families
% SHARD-SUMMARY: Fixes signs and coefficient maps for TFIM, XY, XXZ, Heisenberg, and synthetic grids.
% SHARD-SUMMARY: Distinguishes locally sourced physical labels from synthetic stress-test labels.
% SHARD-SUMMARY: Records the symmetric field split used by every candidate matrix.
% SHARD-KEYWORDS: qubit, Pauli coefficients, TFIM, XY, XXZ, synthetic grid, convention

\section{Pauli Coefficients for Scan Families}
\label{sec:qubit-hamiltonian-family-conventions}

\subsection{Status, Sources, and Dependencies}

\textbf{Status: checked implementation convention.}  The family constructors
are implemented in the Julia backend.

Dependencies: CA-34--CA-53, CONVENTIONS.md (m), (q).

% Source: references/text/GaugingDefectsQuantumSpinSystems.txt:1487--1496.
% Source: references/text/CFTFromLatticeFermions.txt:1351--1403.
% Source: literature/md/2302.14081/2302.14081.md:56--64.
% Checked: test/runtests.jl, testset "qubit Hamiltonian candidate scan".

\subsection{Implemented Maps}

The scan uses \((I,X,Y,Z)\) row/column order.  A one-site field
\(f_a\sigma^a_j\) is represented on bonds by
\[
  h_{a0}=h_{0a}=f_a/2.
\]
Thus the self-dual transverse-Ising sample is
\[
  h=-ZZ-\frac12(XI+IX),
\]
which is the bond-density version of the sourced transverse-Ising form.

The XXZ constructor uses
\[
  h=J(XX+YY+\Delta ZZ),
\]
and Heisenberg is the \(\Delta=1\) specialization.  The transverse-XY
constructor uses \(J_xXX+J_yYY+h_zZ\) with the same symmetric field split.

\subsection{Synthetic Families}

XYZ, compass, DM-style, field-deformed Heisenberg, and deterministic dense
generic matrices are stress tests.  They exercise the gates but carry no
criticality or universality-class claim in this lab book.
